\documentclass[12pt,a4paper]{article}

\usepackage{fontspec}
\usepackage{xunicode}
\usepackage{xltxtra}
\usepackage{polyglossia}
\setmainlanguage{greek}
\setotherlanguage{english}

\setmainfont{DejaVu Serif}
\setsansfont{DejaVu Sans}
\setmonofont{DejaVu Sans Mono}

\usepackage[
top=2.2cm,
bottom=2.2cm,
left=2.2cm,
right=2.2cm
]{geometry}

\usepackage{setspace}
\usepackage{array}
\usepackage{longtable}
\usepackage{booktabs}
\usepackage{titlesec}
\usepackage{enumitem}
\usepackage{xcolor}
\usepackage{hyperref}

\hypersetup{
	colorlinks=true,
	linkcolor=black,
	urlcolor=blue
}

\setstretch{1.15}
\setlength{\parindent}{0pt}
\setlength{\parskip}{0.65em}

\titleformat{\section}
{\large\bfseries}
{}
{0pt}
{}

\begin{document}
	
\begin{center} 
	{\Large \textbf{ΠΡΟΣ ΔΙΚΗΓΟΡΟΥΣ / ΝΟΜΙΚΟΥΣ ΣΥΝΕΡΓΑΤΕΣ}}\\[0.5em] 
	{\large \textbf{Πρόσκληση συνεργασίας γιὰ τὴν ἀξιοποίηση τεκμηριωμένων οἰκονομικῶν καὶ κοινωνικῶν φακέλων σὲ ὑποθέσεις πολιτῶν, ὀφειλετῶν καὶ εὐάλωτων νοικοκυριῶν}} 
\end{center} 

\vspace{1em} 

Ἀξιότιμες κυρίες,\\ 
Ἀξιότιμοι κύριοι, 

Μὲ τὴν παροῦσα ἐπιστολὴ ἐπιθυμοῦμε νὰ σᾶς ἐνημερώσουμε γιὰ μία ὀργανωμένη πρωτοβουλία, ἡ ὁποία ἀφορᾶ τὴν παραγωγὴ τεκμηριωμένων οἰκονομικῶν, κοινωνικῶν καὶ ὑποστηρικτικῶν φακέλων πρὸς ὑποστήριξη πολιτῶν, ὀφειλετῶν, δανειοληπτῶν, ἐπαγγελματιῶν καὶ νοικοκυριῶν ποὺ ἀντιμετωπίζουν πραγματικὴ ἀδυναμία πληρωμῆς ἢ ἀνάγκη θεσμικῆς προστασίας. 

Ἡ πρωτοβουλία αὐτὴ βασίζεται σὲ εἰδικὰ ἐργαλεῖα ἀνάλυσης, μεθοδολογίες ὑπολογισμοῦ τοῦ ἐλάχιστου κόστους διαβίωσης, συστήματα τεκμηρίωσης τῆς πραγματικῆς δυνατότητας πληρωμῆς, καθὼς καὶ σὲ ὀργανωμένη σύνδεση οἰκονομικῶν δεδομένων μὲ πραγματικὰ περιστατικά, ἀποδεικτικὰ στοιχεῖα καὶ τεχνικὲς ἐκθέσεις. 

Σκοπὸς τῆς πρωτοβουλίας δὲν εἶναι ἡ ὑποκατάσταση τοῦ δικηγορικοῦ λειτουργήματος. Ἀντιθέτως, σκοπὸς εἶναι ἡ ἐνίσχυση τοῦ ἔργου τοῦ δικηγόρου μὲ πλήρη, δομημένο καὶ τεκμηριωμένο ὑλικό, τὸ ὁποῖο μπορεῖ νὰ ἀξιοποιηθεῖ, κατὰ τὴν κρίση τοῦ ἁρμόδιου νομικοῦ παραστάτη, σὲ ἐξώδικες ἐνέργειες, δικόγραφα, ἀνακοπές, ἀγωγές, αἰτήσεις ρύθμισης, διοικητικὲς προσφυγὲς καὶ κάθε ἄλλη νόμιμη διαδικασία. 

\section{Αναγκαιότητα συνεργασίας μὲ δικηγόρους} 

Ἡ συνεργασία μὲ δικηγόρους θεωρεῖται ἀπολύτως ἀναγκαία, διότι κάθε ὑπόθεση ποὺ παράγει ἔννομες συνέπειες, κάθε δικόγραφο ποὺ κατατίθεται ἐνώπιον δικαστηρίου ἢ διοικητικῆς ἀρχῆς καὶ κάθε νομικὸς ἰσχυρισμὸς ποὺ προβάλλεται ὑπὲρ πολίτη ἢ ἐπιχείρησης πρέπει νὰ ἀξιολογεῖται, νὰ διατυπώνεται καὶ νὰ ὑπογράφεται ἀπὸ ἁρμόδιο νομικὸ παραστάτη. 

Ἡ τεχνικὴ καὶ οἰκονομικὴ τεκμηρίωση μπορεῖ νὰ ἀναδείξει μὲ ἀκρίβεια τὴν πραγματικὴ κατάσταση τοῦ πολίτη. Δὲν μπορεῖ ὅμως νὰ ἀντικαταστήσει τὴ νομικὴ κρίση, τὴ δικονομικὴ ἐπιλογή, τὴ σύνταξη τοῦ δικογράφου καὶ τὴ δικαστικὴ ἐκπροσώπηση, τὰ ὁποῖα ἀνήκουν ἀποκλειστικὰ στὸν δικηγόρο. 

Γιὰ τὸν λόγο αὐτὸ ἡ συμμετοχὴ δικηγόρων δὲν εἶναι ἁπλῶς χρήσιμη, ἀλλὰ ἀπαραίτητη προϋπόθεση γιὰ τὴν ὀρθή, νόμιμη καὶ ἀποτελεσματικὴ ἀξιοποίηση τῶν μελετῶν καὶ τῶν ἐργαλείων ποὺ ἔχουν ἀναπτυχθεῖ. 

\section{Πεδία ὑποστήριξης} 

Τὰ συστήματα καὶ οἱ μελέτες ποὺ ἔχουν ἀναπτυχθεῖ μποροῦν νὰ προσφέρουν σημαντικὴ ὑποστήριξη στὸ ἔργο τοῦ δικηγόρου, ἰδίως στὰ ἀκόλουθα πεδία: 

\begin{enumerate}[label=\alph*)] 
	\item στὴν πλήρη καταγραφὴ τῆς οἰκονομικῆς κατάστασης τοῦ ἐντολέα, 
	\item στὸν ὑπολογισμὸ τοῦ ἐλάχιστου κόστους διαβίωσης μὲ ἀντικειμενικὰ καὶ ἐπαληθεύσιμα κριτήρια, 
	\item στὴν ἀποτύπωση τῆς πραγματικῆς δυνατότητας πληρωμῆς, 
	\item στὴ συγκέντρωση καὶ ταξινόμηση ἀποδεικτικῶν στοιχείων, 
	\item στὴ δημιουργία τεχνικῶν καὶ οἰκονομικῶν ἐκθέσεων πρὸς χρήση σὲ νόμιμες διαδικασίες, 
	\item στὴ διαμόρφωση φακέλου τεκμηρίωσης γιὰ δικαστικὴ ἢ ἐξώδικη ἀξιοποίηση, 
	\item στὴν ὑποστήριξη ὑποθέσεων ποὺ ἀφοροῦν Δημόσιο, ΑΑΔΕ, ΕΦΚΑ, τράπεζες, servicers, funds, ρυθμίσεις ὀφειλῶν, καταχρηστικὲς πρακτικὲς καὶ ζητήματα πραγματικῆς ἀδυναμίας πληρωμῆς, 
	\item στὴ σύνταξη πινάκων, χρονολογίων, ὑπολογισμῶν καὶ συνοδευτικῶν παραρτημάτων πρὸς ἀξιοποίηση ἀπὸ τὸν δικηγόρο. 
\end{enumerate} 

Ἡ συμβολὴ τοῦ δικηγόρου εἶναι κρίσιμη, διότι μόνο ὁ νομικὸς παραστάτης μπορεῖ νὰ ἀξιολογήσει τὴ δικονομικὴ καταλληλότητα τῶν στοιχείων, νὰ ἐπιλέξει τὸ κατάλληλο ἔνδικο βοήθημα ἢ μέσο, νὰ προσαρμόσει τὴ μελέτη στὶς ἀνάγκες τῆς συγκεκριμένης ὑπόθεσης καὶ νὰ ἀναλάβει τὴν εὐθύνη τῆς τελικῆς νομικῆς διατύπωσης. 

\section{Ρόλος τῆς τεχνικῆς καὶ οἰκονομικῆς ὁμάδας} 

Ἀπὸ τὴν πλευρά μας μποροῦμε νὰ παρέχουμε ὀργανωμένο ὑλικό, τεχνικὲς ἀναλύσεις, οἰκονομικὲς μελέτες, πίνακες, ὑπολογισμούς, τεκμηρίωση, συνοδευτικὲς ἐκθέσεις καὶ προσχέδια δομῆς, τὰ ὁποῖα τίθενται στὴ διάθεση τοῦ συνεργαζόμενου δικηγόρου. 

Ὁ δικηγόρος διατηρεῖ πλήρως τὸν ρόλο του ὡς νομικὸς σύμβουλος, συντάκτης, ἐλεγκτὴς καὶ ὑπεύθυνος χειριστὴς τῆς ὑπόθεσης. Ἡ τεχνικὴ ὁμάδα δὲν παρέχει νομικὲς ὑπηρεσίες, δὲν ὑπογράφει δικόγραφα, δὲν παρέχει δικαστικὴ ἐκπροσώπηση καὶ δὲν ὑποκαθιστᾶ τὴ δικηγορικὴ κρίση. 

Ἡ τεχνικὴ καὶ οἰκονομικὴ ἐργασία λειτουργεῖ ὡς ὑποστηρικτικὸ ἐργαλεῖο, ὥστε ὁ δικηγόρος νὰ ἔχει στὴ διάθεσή του ἕναν ὀργανωμένο φάκελο, μὲ σαφῆ στοιχεῖα, ἀποδείξεις, ὑπολογισμοὺς καὶ συμπεράσματα, τὰ ὁποῖα μπορεῖ νὰ ἀξιοποιήσει κατὰ τὴν ἐπαγγελματική του κρίση. 

\section{Οικονομικό πλαίσιο συνεργασίας} 

Στὸ οἰκονομικὸ σκέλος τῆς συνεργασίας προτείνεται ἡ πλήρης διάκριση μεταξὺ τῆς δικηγορικῆς ἀμοιβῆς καὶ τῆς ἀμοιβῆς τεχνικῆς, οἰκονομικῆς καὶ ὑποστηρικτικῆς τεκμηρίωσης. 

Ἡ ἀμοιβὴ τοῦ δικηγόρου συμφωνεῖται ἀποκλειστικὰ μεταξὺ τοῦ δικηγόρου καὶ τοῦ ἐντολέα του, σύμφωνα μὲ τὸν Κώδικα Δικηγόρων, τὴ δικηγορικὴ δεοντολογία καὶ τὴ φύση τῆς ὑπόθεσης. Ἡ πρωτοβουλία μας δὲν παρεμβαίνει στὸν καθορισμὸ τῆς δικηγορικῆς ἀμοιβῆς, οὔτε ἐπιδιώκει συμμετοχὴ στὴ νομικὴ ἀμοιβὴ τοῦ δικηγόρου. 

Ἀντιστοίχως, ἡ ἀμοιβὴ γιὰ τὴν τεχνική, οἰκονομική, ὑπολογιστικὴ καὶ ἀποδεικτικὴ τεκμηρίωση ἀφορᾶ αὐτοτελῆ ὑποστηρικτικὴ ἐργασία, ὅπως ἡ σύνταξη οἰκονομικῶν μελετῶν, ἡ ὀργάνωση ἀποδεικτικῶν στοιχείων, ἡ παραγωγὴ πινάκων, ἡ ἀνάλυση τῆς πραγματικῆς δυνατότητας πληρωμῆς καὶ ἡ δημιουργία φακέλου τεκμηρίωσης πρὸς ἀξιοποίηση ἀπὸ τὸν δικηγόρο. 

Μὲ τὸν τρόπο αὐτὸ διασφαλίζεται ἡ καθαρότητα τῶν ρόλων, ἡ διαφάνεια πρὸς τὸν πολίτη, ὁ σεβασμὸς στὸ δικηγορικὸ λειτούργημα καὶ ἡ βιωσιμότητα τῆς συνεργασίας. 

\section{Κλιμακωτό μοντέλο πρόσβασης} 

Ἐπειδὴ μεγάλο μέρος τῶν πολιτῶν ποὺ ἀπευθύνονται σὲ τέτοιες διαδικασίες βρίσκεται ἤδη σὲ οἰκονομικὴ ἀδυναμία, τὸ προτεινόμενο μοντέλο δὲν βασίζεται σὲ ὑψηλὸ ἀρχικὸ κόστος. Ἀντιθέτως, προβλέπει χαμηλὴ εἴσοδο, σταδιακὴ ἀνάπτυξη φακέλου καὶ πλήρη τεκμηρίωση μόνο ὅπου ἡ ὑπόθεση τὸ ἀπαιτεῖ. 

Ἐνδεικτικά, τὸ μοντέλο μπορεῖ νὰ ὀργανωθεῖ ὡς ἑξῆς: 
	
	\begin{longtable}{p{0.26\textwidth} p{0.48\textwidth} p{0.18\textwidth}} 
		\toprule 
		\textbf{Επίπεδο} & \textbf{Περιγραφή} & \textbf{Ενδεικτική ἀμοιβή} \\ 
		\midrule 
		Κοινωνικὴ εἴσοδος & Ἄνοιγμα φακέλου, βασικὴ καταγραφή, πρώτη ταξινόμηση στοιχείων καὶ ἔνταξη σὲ ὀργανωμένη διαδικασία. & 20--50€ \\ 
		\addlinespace 
		Βασικὸς φάκελος & Βασικὴ ἀποτύπωση οἰκονομικῆς κατάστασης, ἀρχικὸς ὑπολογισμὸς ἐλάχιστου κόστους διαβίωσης, ἁπλῆ ταξινόμηση ἀποδεικτικῶν. & 120--180€ \\ 
		\addlinespace 
		Κανονικὴ μελέτη & Ἀναλυτικότερη οἰκονομικὴ ἐπεξεργασία, πίνακες, τεκμηρίωση πραγματικῆς δυνατότητας πληρωμῆς καὶ συνοπτικὴ ἔκθεση πρὸς ἀξιοποίηση. & 300--450€ \\ 
		\addlinespace 
		Πλήρης φάκελος τεκμηρίωσης & Πλήρης ἀνάλυση, ἀποδεικτικὸ παράρτημα, ὀργανωμένος φάκελος πρὸς δικηγόρο ἢ δημόσια ἀρχή. & 700--1.350€ \\
		\addlinespace 
		Εἰδικὸς / δικαστικὸς φάκελος & Σύνθετες ὑποθέσεις μὲ μεγάλο οἰκονομικὸ ἀντικείμενο, τραπεζικὴ ζημία, πλειστηριασμούς, ἀποζημιώσεις ἢ ἀγωγές. & ἀπὸ 1.500€ καὶ ἄνω \\ 
		\bottomrule 
	\end{longtable} 
	
Οἱ ἀνωτέρω ἀμοιβὲς ἔχουν ἐνδεικτικὸ χαρακτῆρα καὶ προσαρμόζονται ἀνάλογα μὲ τὴ φύση, τὴν ἔκταση, τὴν πολυπλοκότητα, τὸν ἀριθμὸ τῶν ἐγγράφων, τὴν ἀνάγκη εἰδικῶν ὑπολογισμῶν καὶ τὸν βαθμὸ τεκμηρίωσης ποὺ ἀπαιτεῖ κάθε ὑπόθεση. 

\section{Συλλογικές καὶ θεματικὲς δράσεις} 

Ἰδιαίτερη σημασία ἔχει ἡ δυνατότητα συλλογικῆς ἢ θεματικῆς ὀργάνωσης ὑποθέσεων. Σὲ περιπτώσεις ὅπου πολλοὶ πολῖτες ἀντιμετωπίζουν ὅμοια ἢ παρόμοια προβλήματα, μπορεῖ νὰ δημιουργηθεῖ κοινὴ βάση τεκμηρίωσης, κοινὸ ὑπόδειγμα οἰκονομικῆς ἀνάλυσης, κοινὴ μεθοδολογία καὶ κοινὸς κορμὸς ἐπιχειρημάτων πρὸς ἀξιολόγηση ἀπὸ συνεργαζόμενους δικηγόρους. 

Μὲ τὸν τρόπο αὐτὸ μειώνεται τὸ κόστος γιὰ τὸν πολίτη, αὐξάνεται ἡ ποιότητα τοῦ φακέλου, διευκολύνεται τὸ ἔργο τοῦ δικηγόρου καὶ δημιουργεῖται σοβαρότερη κοινωνικὴ καὶ θεσμικὴ βαρύτητα. 

Ἡ συλλογικὴ ὀργάνωση δὲν ἀναιρεῖ τὴν ἀνάγκη ἐξατομικευμένης νομικῆς ἀξιολόγησης. Κάθε ὑπόθεση παραμένει ἀτομικὴ ὡς πρὸς τὰ πραγματικὰ περιστατικά, τὰ ἀποδεικτικὰ στοιχεῖα καὶ τὴ νομική της πορεία. Ὡστόσο, ἡ κοινὴ τεχνικὴ καὶ οἰκονομικὴ βάση μπορεῖ νὰ περιορίσει κόστος, χρόνο καὶ ἐπαναλαμβανόμενη ἐργασία. 

\section{Οφέλη τῆς συνεργασίας} 

Ἡ συνεργασία αὐτὴ μπορεῖ νὰ λειτουργήσει πρὸς ὄφελος ὅλων τῶν πλευρῶν: 

\begin{itemize} 
	\item Γιὰ τὸν πολίτη, διότι ἀποκτᾶ τεκμηριωμένη εἰκόνα τῆς πραγματικῆς του κατάστασης. 
	\item Γιὰ τὸν δικηγόρο, διότι παραλαμβάνει ὀργανωμένο ὑλικὸ καὶ δὲν ἀναγκάζεται νὰ ξεκινᾶ κάθε ὑπόθεση ἀπὸ μηδενικὴ βάση. 
	\item Γιὰ τὴ δικαιοσύνη, διότι διευκολύνεται ἡ οὐσιαστικὴ ἀξιολόγηση τῆς ὑπόθεσης μὲ συγκεκριμένα δεδομένα. 
	\item Γιὰ τὸ δημόσιο συμφέρον, διότι προωθεῖται μία πιὸ δίκαιη, διαφανὴς καὶ τεκμηριωμένη ἀντιμετώπιση τῆς πραγματικῆς οἰκονομικῆς ἀδυναμίας. 
\end{itemize} 

Στὸ πλαίσιο αὐτό, ἀπευθύνουμε πρόσκληση σὲ δικηγόρους ποὺ ἐνδιαφέρονται νὰ συμμετάσχουν σὲ ἕνα σοβαρὸ δίκτυο συνεργασίας, μὲ σεβασμὸ στὸν θεσμικό τους ρόλο, στὴν ἀνεξαρτησία τους καὶ στοὺς κανόνες τῆς δικηγορικῆς δεοντολογίας. 

\section{Δήλωση ὁρίων καὶ θεσμικοῦ σεβασμοῦ} 

Ἡ παροῦσα πρωτοβουλία δὲν ἀποσκοπεῖ στὴν παροχὴ νομικῶν ὑπηρεσιῶν ἀπὸ μὴ δικηγόρους. Ἀφορᾶ τὴν παραγωγὴ τεκμηριωμένου βοηθητικοῦ ὑλικοῦ, τὸ ὁποῖο μπορεῖ νὰ ἀξιοποιηθεῖ μόνο μέσα ἀπὸ τὴ νόμιμη καὶ ὑπεύθυνη δικηγορικὴ ἐπεξεργασία. 

Κάθε τελικὴ νομικὴ κρίση, κάθε διατύπωση ἰσχυρισμοῦ, κάθε δικονομικὴ ἐπιλογὴ καὶ κάθε δικαστικὴ ἢ ἐξώδικη ἐνέργεια ἀνήκει ἀποκλειστικὰ στὸν δικηγόρο ποὺ ἀναλαμβάνει τὴν ὑπόθεση. 

Ἡ δική μας συμβολὴ περιορίζεται στὴν τεχνική, οἰκονομική, ὀργανωτικὴ καὶ ἀποδεικτικὴ ὑποστήριξη τοῦ φακέλου, ὥστε ὁ δικηγόρος νὰ ἔχει στὴ διάθεσή του πληρέστερο ὑλικὸ καὶ ὁ πολίτης νὰ προσέρχεται στὴ νομικὴ διαδικασία πιὸ ὀργανωμένος καὶ τεκμηριωμένος. 

\section{Επίλογος} 

Ἡ σημερινὴ κοινωνικὴ καὶ οἰκονομικὴ πραγματικότητα ἀπαιτεῖ νέα ἐργαλεῖα τεκμηρίωσης, ἀλλὰ καὶ αὐστηρὸ σεβασμὸ στοὺς θεσμικοὺς ρόλους. Οἱ πολῖτες ποὺ βρίσκονται σὲ πραγματικὴ ἀδυναμία πληρωμῆς δὲν χρειάζονται γενικὲς διακηρύξεις. Χρειάζονται ὀργανωμένο φάκελο, στοιχεῖα, ὑπολογισμούς, ἀποδείξεις καὶ ὑπεύθυνη νομικὴ καθοδήγηση. 

Γιὰ τὸν λόγο αὐτὸ θεωροῦμε ὅτι ἡ συνεργασία μεταξὺ τεχνικῆς/οἰκονομικῆς τεκμηρίωσης καὶ δικηγορικῆς ἐκπροσώπησης μπορεῖ νὰ δημιουργήσει ἕνα οὐσιαστικό, νόμιμο καὶ ἀποτελεσματικὸ πλαίσιο προστασίας πολιτῶν καὶ ὑποστήριξης δικηγόρων. 

Μὲ ἐκτίμηση, 

\vspace{2em} 

\begin{center}
	\begin{minipage}{15cm}
	\textbf{Δημήτριος Μ. Καΐπης, Πολιτικός Μηχανικός ΑΠΘ \& Μηχανικός του Λογισμικού}
\textbf{Τηλέφωνο: 6973751829}\\ 
\textbf{Email: dim.kaipis@gmail.com}\\ 
\textbf{https://m.12com.org} 
	\end{minipage}
\end{center}

\end{document} 
